\chapter{Free Will as Categorical Selection}
\label{ch:free_will}

\begin{chapterabstract}
The free will debate (determinism vs.\ libertarian agency) is reframed within the categorical completion framework. Each conscious moment, the system selects one completion from $10^{6 \times 10^{35}}$ available configurations. This selection is: (i)~physically constrained---by cytoplasmic state, neural connectivity, and BMD history; (ii)~not algorithmically determined---the constraint space is so large that no prior physical state uniquely determines the completion; (iii)~causally efficacious---the selected completion shapes future cytoplasmic states and thus future selections. This is neither classical determinism nor uncaused libertarian agency: it is constrained stochastic selection in categorical space, which is the only form of freedom compatible with physical law and the only form capable of grounding moral responsibility.
\end{chapterabstract}

\section{The Classical Framing}

The free will debate has historically been structured around two options: (1)~determinism---every mental event is the inevitable consequence of prior physical events under causal laws; (2)~libertarian free will---some mental events are not fully determined by prior physical events, implying a causal gap that is either random or mysterious.

Neither option is satisfying. Determinism appears to dissolve the agent; libertarianism appears to dissolve the coherence of agency.

\section{Categorical Selection}

The categorical completion framework provides a third option. At each conscious moment, the selection among $10^{6 \times 10^{35}}$ weak-force configurations is:

\begin{itemize}
  \item \textbf{Not pre-determined}: the constraint space (cytoplasmic state, connectivity, BMD history) narrows the available completions from $10^{6 \times 10^{35}}$ to a much smaller set, but does not determine which element of that set is selected. The selection involves genuine indeterminacy at the level of weak-force angle, dipole orientation, and vibrational phase.

  \item \textbf{Not random}: the narrowing is systematic and sensitive to the agent's history. Different histories produce different constraint sets; the selection is sensitive to who the agent is.

  \item \textbf{Causally efficacious}: the selected completion becomes part of the BMD history and shapes future constraint sets. Choices have consequences.
\end{itemize}

\section{Implications for Moral Responsibility}

Moral responsibility requires that agents be the appropriate sources of their actions. Categorical selection satisfies this: actions originate in the agent's constraint set (their character, history, knowledge) and are sensitive to reasons (which alter the cytoplasmic state and thereby the constraints). This is compatible with physical law and grounds moral responsibility without requiring causal gaps.
